\documentclass{article}

\usepackage{amsmath,amssymb}
\usepackage{xurl}
\usepackage[hidelinks]{hyperref}
\setlength{\emergencystretch}{2em}

\input{command}

\title{Scalar Normalization in Wolf's Lorentz Normal Form}
\author{}
\date{2026-07-16}

\begin{document}
\maketitle
\gapnote{open-gap}{open}

\section{The assertion}

Wolf's Proposition~2.11 states that every qubit channel
$T:\MN{2}\to\MN{2}$ admits invertible completely positive maps
$\Phi_1,\Phi_2$, each of Kraus rank one, such that
\begin{align}
  T' = \Phi_2\circ T\circ\Phi_1
\end{align}
is again a qubit channel and belongs to one of three Lorentz normal forms:
diagonal, non-diagonal, or singular.\footnote{\cite[Proposition~2.11]{Wolf2012Quantum};
local source \path{Notes/WolfNoteTexSource/ch02_representations.tex},
lines 1021--1035.}  Each filtering has the form
\begin{align}
  \Phi_X(A)=XAX^\dagger,
  \qquad X\in\mathrm{GL}(2,\C).
\end{align}
Thus the proposition allows every invertible Kraus-rank-one completely positive
filter, not only those represented by determinant-one matrices.

\section{Determinant-one action and scalar freedom}

Choose a square root $c^2=\det X$ and write $X=cS$, where
$S\in\mathrm{SL}(2,\C)$. Then
\begin{align}
  \Phi_X(A)
    &= |c|^2\Phi_S(A),\\
  |c|^2
    &= |\det X|.
\end{align}
Consequently, if $X_i=c_iS_i$ for the two filters, then
\begin{align}
  \Phi_{X_2}\circ T\circ\Phi_{X_1}
  = |c_1c_2|^2
    (\Phi_{S_2}\circ T\circ\Phi_{S_1}).
  \label{eq:scalar-factor}
\end{align}
The matrices $S_1,S_2$ determine the determinant-one Lorentz action. The
positive factor $|c_1c_2|^2$ selects the normalization on the resulting ray of
completely positive maps.

This distinction is immaterial when only the Lorentz orbit up to positive scale
is considered. It is essential when the representative is required to be a
channel. If the determinant-one representative satisfies
\begin{align}
  \tr\!\left(\Phi_{S_2}\circ T\circ\Phi_{S_1}(A)\right)
  = \alpha\,\tr(A),
\end{align}
then it is trace preserving only when $\alpha=1$. General invertible filters
can multiply the map by $\alpha^{-1}$ through the scalar in
\eqref{eq:scalar-factor}; determinant-one filters have no such freedom. The
minimisation in Propositions~2.8 and~2.9 produces partial traces proportional to
the identity, but the proportionality scalar must still be chosen so that the
qubit representative is trace preserving.

\section{The point at issue}

A former local statement required both filters to be determinant-one while also
requiring the filtered map to be a qubit channel in one of the three canonical
forms.\footnote{The removed declaration was
\leanid{Wolf.exists_lorentz_normal_form_qubit} in
\doclink{QICLean/Channel/LorentzNormalForm}.}  Equation~\eqref{eq:scalar-factor}
shows why the determinant-one restriction does not justify the normalization in
Wolf's proposition. The source theorem permits general invertible filters, and
its channel conclusion uses their scalar freedom. The restricted local
statement was therefore removed rather than treated as a formalization of
Proposition~2.11.

The three canonical-form predicates and the square full-Kraus-rank minimisation
theorem are formalized.  The Pauli--Minkowski identification, the
$\mathrm{SL}(2,\C)$ congruence action, its image in
$\mathrm{SO}^+(1,3)$, and the exact transfer-matrix action
$L_2\widehat T L_1$ are also formalized.\footnote{Formal declarations
\leanid{Wolf.IsLorentzDiagonal}, \leanid{Wolf.IsLorentzNonDiagonal},
\leanid{Wolf.IsLorentzSingular}, and
\leanid{Wolf.exists_normal_form_generic} in
\doclink{QICLean/Channel/LorentzNormalForm}, together with
\leanid{Wolf.pauliMinkowskiEquiv},
\leanid{Wolf.spinorMatrix_isSpecialOrthochronousLorentz}, and
\leanid{Wolf.pauliTransferMatrixReal_slFiltering} in
\doclink{QICLean/Channel/LorentzNormalForm/SpinorAction};
\leanid{Wolf.spinorCoverHom_surjective},
\leanid{Wolf.spinorCoverHom_eq_iff_eq_or_eq_neg}, and
\leanid{Wolf.boostSpinor_posDef} in
\doclink{QICLean/Channel/LorentzNormalForm/SpinorCover}; and
\leanid{Wolf.spinorMatrix_boostExpSL2} and
\leanid{Wolf.spinorMatrix_rotationExpSL2} in
\doclink{QICLean/Channel/LorentzNormalForm/SpinorExponential}.}  The theorem
asserting existence of a canonical representative for every qubit channel
remains open.

All three displayed representatives are now formalized directly from the
source formulas in
\cite[Theorem~8, Eqs.~(17)--(19)]{VerstraeteVerschelde2002QuantumChannels}.
Verstraete--Verschelde define $R_\Phi$ from the first-factor-partially-transposed
dual state and then use $(1,x')^{\mathsf T}=R_\Phi(1,x)^{\mathsf T}$, so
$R_\Phi$ is their Bloch transfer matrix.  It corresponds to the local
$\widehat T$, not to the raw Pauli correlation matrix of $\tau$.  In the
normalized Choi convention and the local factor order, the latter satisfies
$R_{\mathrm{raw}}(\tau)=\widehat T\operatorname{diag}(1,1,-1,1)$ because
$\sigma_2^{\mathsf T}=-\sigma_2$.
The partial-transpose sign is therefore already absorbed before the parameters
in Theorem~8 are introduced.

For the diagonal representative, put $q_i=p_i^2$ in the Equation~(18) Pauli
family and specialize to $A=B=\one$.  Direct Pauli conjugation gives
\begin{align}
  s_1&=q_0+q_1-q_2-q_3,
  &s_2&=q_0-q_1+q_2-q_3,\\
  s_3&=q_0-q_1-q_2+q_3,
  &1&=q_0+q_1+q_2+q_3.
\end{align}
Inverting these equations in Wolf's Pauli-transfer convention gives
\begin{align}
  q_0&=(1+s_1+s_2+s_3)/4,
  &q_1&=(1+s_1-s_2-s_3)/4,\\
  q_2&=(1-s_1+s_2-s_3)/4,
  &q_3&=(1-s_1-s_2+s_3)/4.
\end{align}
Verstraete--Verschelde, Theorem~8 instead prints
$1-s_1-s_2-s_3\geq0$.  This is inconsistent with Equation~(16): the identity
channel has $R_\Phi=\operatorname{diag}(1,1,1,1)$ and violates that printed
inequality.  It also cannot be explained by the partial-transpose bridge,
because $R_\Phi$ is defined only after that transpose and already acts as the
Bloch transfer matrix.  The formalization therefore records the all-minus
condition as a printed sign inconsistency and does not use it.
The formalization proves that the weights sum to one and that their simultaneous
nonnegativity is equivalent to the four displayed numerator inequalities.
Assuming that nonnegativity, the Equation~(18) amplitudes are $p_i=\sqrt{q_i}$
when $A=B=\one$, and the resulting channel has Pauli transfer matrix
$\operatorname{diag}(1,s_1,s_2,s_3)$ and Choi/Kraus rank
$\#\{i\mid q_i\ne0\}$, retaining all rank-drop boundaries.  No necessity
statement for complete positivity is formalized: the definition uses
$\sqrt{q_i}$ and therefore clips a negative candidate weight to zero.
Under $1\geq s_1\geq s_2\geq|s_3|$, Wolf's Equation~(2.40),
$s_1+s_2\leq1+s_3$, supplies all four inequalities and agrees with the direct
Equation~(18) calculation above.
The exact non-diagonal three-Kraus family defines a channel for each assumed
$x\in[0,1]$, has the Pauli data
$\Delta=\operatorname{diag}(x/\sqrt3,x/\sqrt3,1/3)$ and
$v=(0,0,2/3)$, and has Choi/Kraus rank three for $x<1$ and two at $x=1$.
The third normal form in Equation~(17) is realized by the singular two-Kraus
family
$X\mapsto\tr(X)\lvert0\rangle\!\langle0\rvert$ and has Choi/Kraus rank two.
The diagonal rank formula is the direct Bell-family calculation from
Equation~(18).  The non-diagonal and singular rank statements are respectively
cases 2 and 3 of \cite[Theorem~18]{WolfCirac2008Dividing}\footnote{Formal declarations
\leanid{Wolf.diagonalMap},
\leanid{Wolf.diagonalBellWeight_nonneg_iff},
\leanid{Wolf.pauliTransferMatrix_diagonalMap},
\leanid{Wolf.choiRank_diagonalMap},
\leanid{Wolf.nonDiagonalMap},
\leanid{Wolf.choiRank_nonDiagonalMap_eq_three},
\leanid{Wolf.choiRank_nonDiagonalMap_one},
\leanid{Wolf.singularMap}, and
\leanid{Wolf.choiRank_singularMap} in
\doclink{QICLean/Channel/LorentzNormalForm/CanonicalQubitChannels}.}.
This bounded construction proves sufficiency for an already supplied
$x\in[0,1]$; it does not derive that range or classify Lorentz orbits.

\section{Elimination plan}

A source-faithful proof requires the following steps.
\begin{enumerate}
  \item[\(\checkmark\)] General invertible Kraus-rank-one completely positive
    filters are now formalized as \leanid{Wolf.InvertibleFilter}.  For every
    nonzero dimension,
    \leanid{Wolf.InvertibleFilter.exists_scalar_slFiltering} separates the
    filtering matrix as $X=cS$, where $c\in\C$ is nonzero,
    $c^d=\det X$, and $S\in\operatorname{SL}(d,\C)$, while the induced map satisfies
    $\Phi_X=|c|^2\Phi_S$ with $|c|^2>0$.  The corresponding two-filter scalar
    identity for $\Phi_2\circ T\circ\Phi_1$ is
    \leanid{Wolf.InvertibleFilter.filteredMap_eq_normSq_smul}.
  \item Use the minimisation argument of Propositions~2.8 and~2.9 for the
    full-Kraus-rank case, retaining the proportionality constants of the two
    marginals.
  \item[\(\checkmark\)] The Hermitian congruence $M\mapsto SMS^\dagger$
    for $S\in\mathrm{SL}(2,\C)$ is now identified with a real Minkowski
    transformation \leanid{Wolf.spinorMatrix}.  Its image satisfies the three
    defining conditions of $\mathrm{SO}^+(1,3)$, and two-sided filtering acts
    on Pauli transfer matrices in Wolf's order as $L_2\widehat T L_1$.
  \item[\(\checkmark\)] The spinor homomorphism onto
    $\mathrm{SO}^+(1,3)$ is surjective and has the exact two-point fibres
    $\{S,-S\}$.  The explicit boost and rotation exponential identities from
    Wolf's Equation~(2.44) are also formalized, so arbitrary source Lorentz
    factors can be lifted by the proved spinor cover rather than extrapolating
    from the $\mathrm{SU}(2)$ rotation case.
  \item Classify the resulting diagonal, non-diagonal, and singular Lorentz
    orbits using the source decomposition and its signed canonical-pair data.
  \item Choose the scalar factors so that the canonical representative is trace
    preserving, and derive the stated parameter bounds from the positive
    two-qubit normal forms.  Once the Bell-diagonal eigenvalues are nonnegative and
    $x\in[0,1]$ is supplied, the exact Kraus constructions, complete positivity,
    transfer data, and rank calculations for all three representatives are
    already formalized.
  \item Prove the three-case existence theorem using the general filters and link
    it to the existing canonical-form predicates.
\end{enumerate}

The discrepancy is an open gap: the canonical predicates, the exact
displayed representatives and their ranks, the generic minimisation theorem
in square dimension, the determinant-one Lorentz action, and the full spinor
cover are available, but the scalar-normalized Lorentz-orbit classification
and necessity of the parameter range required by Proposition~2.11 have not
yet been formalized.

\bibliographystyle{alpha}
\bibliography{references}

\end{document}
